% r03_h1s_d: proof of the two-flow comparison bound, part 4: first moments of the units, the vacuum polarisation, the slice law, gradients of externals, the compound as a weakened element.

\begin{lemma}[the first moments of the units vanish]\label{rot:h1s:M}
In each flow's own frame: \textup{(i)} for every slice assignment $\sigma$ the self-energy kernel $k_\sigma$, as it enters any family (the sum of the two Wick contractions between the same two legs, in leg-site coordinates), has $F_\sigma\defeq\sum_uk_\sigma(u)\,u=0$; \textup{(ii)} for every unordered slice assignment $\sigma=\{j_1,j_2\}$ of the vacuum polarisation, the line-symmetrised kernel $\Pi_{\sigma,\mu\nu}\defeq\Pi_{(j_1,j_2),\mu\nu}+\Pi_{(j_2,j_1),\mu\nu}$ has $F_{\sigma,\mu\nu\rho}\defeq\sum_y\Pi_{\sigma,\mu\nu}(x,y)(y-x)_\rho=0$, in site and in bond-midpoint coordinates alike. (For a labelled assignment $(j_1,j_2)$ with $j_1\ne j_2$ the first moment is not zero; this is why (NF4) merges the two classes.)
\end{lemma}
\begin{proof}
Why it is needed: the compound's expansion $\sum_yk_\sigma(x,y)f(y)=c_\sigma f(x)+F_\sigma\cdot\nabla f(x)+\sum_yk_\sigma(x,y)T_xf(y)$ removes the first-order term. For a slice-line leg a term $F_\sigma\cdot\nabla f(x)$ would be harmless even if $F_\sigma\ne0$ ($F^Y=RF^X$ by \Cref{rot:h1s:Z0}). For an external or coarse leg it would not: $F^X\cdot\nabla E^X(x)$ and $F^X\cdot R^{-1}\nabla E^Y(\tau x)$ differ by $F^X\cdot(1-R^{-1})\nabla E(x)$ up to small terms, a frame term of order $|F|\|E\|$ with no gain in $\eta'$. (i) After (NF2) the self-energy at slice $\sigma$ is a bilinear form $B_\sigma(\phi,\phi)=\sum_{x',y'}\phi(x')\cdot\tilde k_\sigma(x',y')\phi(y')$ in the two leg fields, $x'$, $y'$ the leg sites (vertex sites shifted by the bonds of the currents), with $\tilde k_\sigma$ proportional to the identity in the $O(2)$ indices ($q^2=-1$; no $q$-part by charge conjugation). In every family the self-energy sits between two leg elements $g$ at $x'$ and $f$ at $y'$ through both Wick contractions, which have identical element sets; their sum is $\sum g(x')[\tilde k_\sigma(x',y')+\tilde k_\sigma(y',x')]f(y')$. The kernel $k_\sigma\defeq\tilde k_\sigma+\tilde k_\sigma^T$ is translation invariant and symmetric, so $\kappa_\sigma(u)\defeq k_\sigma(x',x'+u)$ is even and $F_\sigma=0$, whatever the rewriting of \Cref{rot:h1s:cv} did inside the kernel; the expansion is about the leg site, so no shifted-site moment appears. (ii) The bubble at labelled slices is $\Pi_{(j_1,j_2),\mu\nu}(x,y)\propto C^{(j_1)}(x+e_\mu,y+e_\nu)C^{(j_2)}(x,y)-C^{(j_1)}(x+e_\mu,y)C^{(j_2)}(x,y+e_\nu)$; under $-\mathbf 1\in\Oh$ the two factor positions swap slices, so the labelled kernel is $\Oh$-covariant only up to $j_1\leftrightarrow j_2$ and the symmetrised one is the $\Oh$-covariant object, a kernel on bonds covariant in midpoint coordinates, $\Pi_\sigma(Rm,Rm')=R\Pi_\sigma(m,m')R^T$ for $R\in\Oh$. Its midpoint first moment is an $\Oh$-invariant rank-three tensor, hence $0$; its zero moment is $c_\sigma\delta_{\mu\nu}$. In site coordinates $m'-m=(y-x)+(e_\nu-e_\mu)\eta'/2$, so $F_{\sigma,\mu\nu\rho}=0-c_\sigma\delta_{\mu\nu}(e_\nu-e_\mu)_\rho\eta'/2=0$.
\end{proof}

\begin{lemma}[the vacuum polarisation]\label{rot:h1s:VP}
\textup{(a)} With every line sliced, the order-$e^2$ photon two-point function at zero momentum vanishes in each flow, \cite[(2.23)--(2.28)]{Ba3} (valid with free boundary conditions) and \cite[\S3.5]{K1}: $\sum_{\sigma}c_{\sigma,\mu\nu}+\sum_jc^{\rm sg}_j\delta_{\mu\nu}+c^{\rm dt}_{\mu\nu}+c^{\rm top}_{\mu\nu}=0$, where $c_\sigma\defeq\sum_y\Pi_\sigma(x,y)$ is the bubble's zero moment, $c^{\rm sg}_j$ the seagull tadpole, $c^{\rm dt}$ the derivative tadpole and $c^{\rm top}$ the top pieces, all common numbers, all sums finite ($\sum_\sigma|c_\sigma|\le C\eta'^{-1}$ at fixed $r$). \textup{(b)} The seagull and derivative tadpoles are single-vertex structures of degree $-1$ and may not appear as elements; (a) replaces them by $-\sum_\sigma c_\sigma-c^{\rm top}$. With the local terms $+c_\sigma f(x)g(x)$ of the bundled bubbles ((R1); $f$, $g$ the photon legs) the complete local content at a vacuum-polarisation vertex is $[-\sum_{\sigma\ \rm unbundled}c_\sigma-c^{\rm top}]f(x)g(x)$, a slice assignment being unbundled iff $\sigma_{\max}>J\defeq\max(j_f,j_g)$, plus the remainders of the bundled bubbles and the unbundled bubbles themselves, which are never alone ($\{f,\Pi_\sigma\}$ has $D=1$). Since $|c_\sigma|\le C\ell_{\sigma_{\min}}/\ell_{\sigma_{\max}}^2$ (placement $(1,1)$) and $\le C\ell_{\sigma_{\min}}^2/\ell_{\sigma_{\max}}^3$ (placement $(2,0)$, using $\sum_u\partial\partial C^{(j_1)}(u)=0$),
\[
\Big|\sum_{\sigma_{\max}>J}c_\sigma\Big|\le C\sum_{j_2>J}\ell_{j_2}^{-2}\sum_{j_1\le j_2}\ell_{j_1}\le C\ell_J^{-1},\qquad|c^{\rm top}|\le C:
\]
a vertex weight $-1$ at scale $J\ge j(x)$ plus a weight $0$. Photon legs carry no own derivative, so $\{f,\text{weight}\}$ has $D=1$ and $\{f,g,\text{weight}\}$ has $D=3$. The Ward identity is used once, per flow, for the sum over all $\sigma$; its fine end, $\sum_{\sigma_{\max}\le J}c_\sigma$, of order $\eta'^{-1}$, is exactly what the bundled remainders carry, and only the coarse end survives as a weight. With external photon legs the leftover is $0$ exactly. \textup{(c)} The order-$e^4$ vacuum polarisation, and every degree-zero two-point structure handled by (R2b), needs no identity: $c_{\mu\nu}\defeq\sum_t\Pi^{(4)}_W(x,t)$ is a common number (\Cref{rot:h1s:Z0}), isotropic, and $|\sum_\sigma c_\sigma|\le C(1+r)^{n_0+1}$ by \Cref{rot:h1s:MS} applied per flow to W with one vertex fixed (all proper components of W are already of degree $\ge0$ after the rule; $D(W)=0$ gives one logarithm). The local term $c\,A^F(x)\cdot A^F(x)$ is a weight-$0$ vertex term common to both flows; its difference is telescoped over the externals. \textup{(d)} In the scalar sector the bundled self-energy's zeroth-order term cancels against ${\rm ct}_\sigma$ exactly; the sunset and the scalar $e^4$ structures give $c=0$ exactly with the co-switched order-four share.
\end{lemma}

\begin{lemma}[the slice law]\label{rot:h1s:P}
Let $g_j\defeq C^{(j)}-C^{(j)}_c$ with $C^{(j)}_c(v)\defeq\int_{\ell_j^2}^{\ell_{j+1}^2}e^{-m_r^2s}(4\pi s)^{-3/2}e^{-|v|^2/4s}ds$ the isotropic continuum slice, $3\le j\le r-1$. Then for every $n\ge0$ lattice derivatives
\[
|\partial^ng_j(u)|\le C_n(\eta'/\ell_j)^{1/2^n}\ell_j^{-1-n}e^{-\delta_n|u|/\ell_j},
\]
and consequently, for matched points $|u'-R^{-1}u|\le4c_\tau\eta'$ (bond shifts included),
\[
|\partial^nC^{(j)}(u)-R^{\otimes n}\partial^nC^{(j)}(u')|\le C(\eta'/\ell_j)^{1/2^n}\ell_j^{-1-n}e^{-\delta_n|u|/(2\ell_j)},
\]
the same for the photon slices and, with $\ell_j$ replaced by $1$ and the mass decay, for the top piece. For $j\le2$ there is no gain and the factor $(\eta'/\ell_j)^\theta$ is paid by a degree loss $\theta$ (\Cref{rot:cor:unif}(a), King's device at the finest slices \cite[p.~675]{K1}). Constants are uniform in $m_r\in(0,1]$ and $r$.
\end{lemma}
\begin{proof}
Value clause $n=0$: the coordinates of the continuous-time walk are independent one-dimensional walks \cite[Prop.~1.2.2]{LL10}, and the local limit law \cite[Thm.~2.5.6]{LL10} gives, for $|u|\le s/(8C_0\eta')$ and $s\ge\ell_3^2$, $|\eta'^{-3}P_s(u)-(4\pi s)^{-3/2}e^{-|u|^2/4s}|\le c(\eta'/\sqrt s)(4\pi s)^{-3/2}e^{-|u|^2/8s}$; integrated over $s\in[\ell_j^2,\ell_{j+1}^2]$ this is $\le c(\eta'/\ell_j)\ell_j^{-1}e^{-\delta|u|/\ell_j}$, and beyond that range both kernels are bounded by their Gaussian and Poisson tails, $\le C_N(\eta'/\ell_j)^Ne^{-\delta|u|/\ell_j}$; the geometric mean with the decay of each kernel halves the rate and keeps the factor. Derivatives by interpolation: for a lattice function $h$ and $\rho=N\eta'$, $\partial_\mu h(u)=[h(u+\rho e_\mu)-h(u)]/\rho-\rho^{-1}\sum_{i=1}^{N-1}(\rho-i\eta')\partial_\mu\partial_\mu h(u+i\eta'e_\mu)\eta'$, so $|\partial_\mu h(u)|\le2\sup|h|/\rho+\rho\sup|\partial\partial h|$ over the segment. With $h=\partial^ng_j$, the crude bound $|\partial^{n+2}g_j|\le|\partial^{n+2}C^{(j)}|+|\partial^{n+2}C_c^{(j)}|\le C\ell_j^{-3-n}e^{-\delta|u|/\ell_j}$ from \eqref{eq:h1s:B} and the smoothness of $C_c$, and $\rho\defeq\ell_j\epsilon_n^{1/2}$ with $\epsilon_n\defeq C_n(\eta'/\ell_j)^{1/2^n}$ (so $\rho\ge\eta'$; round to a multiple of $\eta'$), one gets $|\partial^{n+1}g_j|\le C\epsilon_n^{1/2}\ell_j^{-2-n}$, the decay rate kept since the segment has length $\le\ell_j$; induction on $n$. Matched points: $\partial^nC^{(j)}(u)-R^{\otimes n}\partial^nC^{(j)}(u')=\partial^ng_j(u)+[\partial^nC_c^{(j)}(u)-R^{\otimes n}\partial^nC^{(j)}_c(u')]-R^{\otimes n}\partial^ng_j(u')$; the middle term is the rotation covariance of the isotropic $C_c$, $\nabla^nC_c(Rv)=R^{\otimes n}\nabla^nC_c(v)$, plus the displacement $|R^{-1}u-u'|\le4c_\tau\eta'$ against $|\nabla^{n+1}C_c|\le C\ell_j^{-2-n}$, and the lattice-difference-versus-derivative error of $C_c$, $\le C\eta'\ell_j^{-2-n}$; both are inside the claim. (The true law is $(\eta'/\ell_j)^2$ for the values and $\eta'/\ell_j$ for first differences; the exponents $1/2^n$ are far from sharp and suffice, since every gain is compared with $\theta_1\le\frac1{16}$.)
\end{proof}

\begin{lemma}[gradients of externals and coarse elements]\label{rot:h1s:E}
For $E\in\{H^F\phi_r,H^F_AA_r$, their covariant derivatives, the background$\}$ and for $\Gamma^F(\cdot,z)$, $H^F(\cdot,z)$: $|E^X(x)-E^Y(\tau x)|\le C\eta^{1/2}\|E\|$ (\Cref{rot:lem:H}, \Cref{rot:lem:Gam}(iii)), and both are restrictions of $C^{1,\alpha}$ functions $g^X$, $g^Y$ on $\R^3$ with $\|\nabla g^F\|_\infty\le C_\nabla\|E\|$ and $[\nabla g^F]_\alpha\le C\|E\|$ (\cite[Thm.~3.3, (3.8)]{K1} for the standard flow, \Cref{rot:prop:BD}(a) for the hybrid flow; for the lifts, the Fourier representation of \Cref{rot:lem:H} with smooth symbols). Then, with $\kappa\defeq\alpha/(2(1+\alpha))$, $|\nabla g^X(x)-\nabla g^Y(x)|\le C\eta^\kappa\|E\|$ for all $x\in\R^3$, in particular at $x$ and at $\tau x$. Consequently the second-order own-frame Taylor remainders $T^F_xE^F(y)\defeq E^F(y)-E^F(x)-(y-x)\cdot\nabla g^F(x)$ (in the $Y$ frame at $\tau x$, $\tau y$) differ by
\[
|T^X_xE^X(y)-T^Y_{\tau x}E^Y(\tau y)|\le C\|E\|\min\{|u|^{1+\alpha}+\eta'^{1+\alpha},\ \eta^{1/2}+|u|\eta^\kappa+\eta'\},\qquad u=y-x .
\]
For a fine slice line $f=\partial^aC^{(j_f)}(\cdot-z)$ in place of $E$ the same holds with $\eta^{1/2}$, $\eta^\kappa$, $\eta'$ replaced by $(\eta'/\ell_{j_f})^{1/2^a}\ell_{j_f}^{-1-a}$, $(\eta'/\ell_{j_f})^{1/2^{a+1}}\ell_{j_f}^{-2-a}$, $\eta'\ell_{j_f}^{-2-a}$, times $e^{-\delta\,{\rm dist}/(2\ell_{j_f})}$, and with the remainder bound $c|u|^{1+\alpha}\ell_{j_f}^{-2-a-\alpha}$ (\Cref{rot:h1s:P}; no interpolation needed).
\end{lemma}
\begin{proof}
$h\defeq g^X-g^Y$ has $\|h\|_\infty\le C\eta^{1/2}\|E\|$ (at lattice points by the value difference and $|g^Y(\tau x)-g^Y(x)|\le C_\nabla c_\tau\eta'\|E\|$, between them by the Lipschitz bound) and $[\nabla h]_\alpha\le2C\|E\|$; for a unit vector $e$ and $\rho>0$, $|e\cdot\nabla h(x)|\le|h(x+\rho e)-h(x)|/\rho+[\nabla h]_\alpha\rho^\alpha\le2C\eta^{1/2}\|E\|/\rho+2C\|E\|\rho^\alpha$; choose $\rho=\eta^{1/(2(1+\alpha))}$. The first remainder bound is the sum of the two remainders ($|\tau y-\tau x|\le|u|+2c_\tau\eta'$); the second is the value law at $y$ and $x$, the gradient law at $x$ and the displacement $|(\tau y-\tau x)-(y-x)|\le2c_\tau\eta'$ against $\|\nabla g^Y\|$.
\end{proof}

\begin{lemma}[the compound as a weakened element]\label{rot:h1s:U}
Let $B=(U,f)$ be a compound of (R1): $U$ a self-energy or vacuum polarisation with slice assignment $\sigma$ between the leg sites $x$ and $y$, placement-summed and symmetrised kernel $k_\sigma(x,y)=\kappa_\sigma(x-y)$ with $|\kappa_\sigma(u)|\le C\ell_{j_1}^{d_1}\ell_{j_2}^{d_2}e^{-\delta|u|/\ell_{\sigma_{\min}}}$, $d_1+d_2=-4$ (the product of two slices decays at the finest line); $f=\partial^aC^{(j_f)}(y-w)$, $a\in\{0,1,2\}$ ($a\le1$ for a compound as formed, \Cref{rot:h1s:nd}(ii); $a=2$ when the compound is the differenced factor of (R2b), see (iv)), $\sigma_{\max}\le j_f$; $K^F_\sigma$ the own-frame second-order remainder sums $K^X_\sigma(x,w)\defeq\sum_{y}\kappa_\sigma(x-y)T^X_xf^X(y)$, $K^Y_\sigma(\tau x,\tau w)\defeq\sum_{y'\in\Lambda_Y}\kappa_\sigma(R^{-1}(\tau x-y'))T^Y_{\tau x}f^Y(y')$, with lattice differences in the Taylor polynomials. Put $s\defeq\alpha/(1+\alpha)$, $\kappa\defeq\alpha/(2(1+\alpha))$.
\begin{enumerate}[label=\textup{(\roman*)},nosep,leftmargin=*]
\item $|\sum_{\sigma:\sigma_{\max}\le j_f}K_\sigma(x,w)|\le C\ell_{j_f}^{-2-a}e^{-\delta|x-w|/(4\ell_{j_f})}$ per flow: the compound is a line of dimension $-2-a$ at scale $j_f$. With an external or coarse leg $E$: $|\sum_{\sigma_{\max}\le j}K_\sigma(x)|\le C\ell_j^\alpha\|E\|$, a vertex weight $\alpha$ at scale $\sigma_{\max}$, summed geometrically.
\item $\sum_{\sigma_{\max}\le j_f}|K^X_\sigma(x,w)-K^Y_\sigma(\tau x,\tau w)|\le C(1+r)\eta'^{\theta_U}\ell_{j_f}^{-2-a-\theta_U}e^{-\delta|x-w|/(4\ell_{j_f})}$ with $\theta_U\defeq\min\{\theta,s/2^{a+1}\}$; with an external or coarse leg $\le C(1+r)\eta'^{\theta_U}\|E\|$ with $\theta_U\defeq\min\{\theta,\kappa s\}$. No frame vector appears.
\item For the double compound $(U_1,l,U_2)$ with $l=\partial^aC^{(j_l)}$, $a\in\{0,1\}$: per flow a line of dimension $-3-a$ at scale $j_l$; the difference gains $(\eta'/\ell_{j_l})^{s/8}$, with one logarithm per unit. The same holds when $l$ is a coarse element $\Gamma^F$: per flow $|\sum_{\sigma_1,\sigma_2}(U_1,\Gamma,U_2)(x,w)|\le Ce^{-m_r|x-w|/2}$ (dimension $0$ at the top scale), and the difference is $\le C(1+r)^2\eta^{\kappa s}e^{-m_r|x-w|/4}$.
\item First-order forms (the compounds $(T,f)$, $(SC,f)$ and the brackets of (R2b) whose differenced factor is a line carrying $a\le2$ derivatives): gain $(\eta'/\ell_{j_f})^{1/2^{a+2}}\ge(\eta'/\ell_{j_f})^{1/16}$ (slice leg) resp.\ $\eta^{1/4}$ (external leg), with one logarithm. When the differenced factor of (R2b) is a compound $(U,f)$, the bracket is bounded by (i)--(ii) with $a+1\le2$ in place of $a$: per flow $|t-x|$ times a line of dimension $-3-a$ at scale $j_f$, and the gain $\eta'^{\theta_U}$ with $\theta_U\ge\min\{\theta,s/8\}$.
\end{enumerate}
\end{lemma}
\begin{proof}
(i) $\sum_u\eta'^3|\kappa_\sigma(u)||u|^{1+\alpha}\le C\ell_{j_1}^{d_1}\ell_{j_2}^{d_2}\ell_{\sigma_{\min}}^{4+\alpha}\le C\ell_{\sigma_{\max}}^\alpha$ (the sum is dominated by $|u|\sim\ell_{\sigma_{\min}}$, where the coarser line is bounded by $\ell_{\sigma_{\max}}^{d_{\max}}$; more precisely it is $\le C\ell_{\sigma_{\min}}^{\alpha}(\ell_{\sigma_{\min}}/\ell_{\sigma_{\max}})^{|d_{\max}|}$, geometric in $\sigma_{\min}$ at fixed $\sigma_{\max}$); $|T_xf(y)|\le C\ell_{j_f}^{-2-a-\alpha}|u|^{1+\alpha}$ with the leg's decay $e^{-\delta|y-w|/(2\ell_{j_f})}\le e^{-\delta|x-w|/(4\ell_{j_f})}e^{\delta|u|/(4\ell_{j_f})}$, the last factor absorbed by the unit's decay at $\ell_{\sigma_{\min}}\le\ell_{j_f}$; $\sum_{\sigma_{\max}\le j_f}\sum_{\sigma_{\min}\le\sigma_{\max}}\ell_{\sigma_{\min}}^{\alpha}(\ell_{\sigma_{\min}}/\ell_{\sigma_{\max}})^{|d_{\max}|}\le C\ell_{j_f}^\alpha$. External leg: the same with $\|E\|$. (ii) Split $K^X_\sigma-K^Y_\sigma=\sum_y[\kappa_\sigma(x-y)-\kappa_\sigma(R^{-1}(\tau x-\tau y))]T^X_xf^X(y)+\sum_y\kappa_\sigma(R^{-1}(\tau x-\tau y))[T^X_xf^X(y)-T^Y_{\tau x}f^Y(\tau y)]=:A_\sigma+B_\sigma$. $A_\sigma$: telescoping the product of the unit's lines with \Cref{rot:h1s:P} ($n\le2$) for the differenced line and \eqref{eq:h1s:B} for the other gives $|\Delta\kappa_\sigma|\le C(\eta'/\ell_{\sigma_{\min}})^{1/4}\ell_{j_1}^{d_1}\ell_{j_2}^{d_2}e^{-\delta|u|/(2\ell_{\sigma_{\min}})}$ (with the $125^\theta$-form for $\sigma_{\min}\le2$), so as in (i) $\sum_{\sigma_{\max}\le j_f}A_\sigma\le C\eta'^\theta\ell_{j_f}^{\alpha-\theta}\ell_{j_f}^{-2-a-\alpha}$ for $\theta\le\frac14$, $\theta<\alpha$. $B_\sigma$: the bracket is $\le\min\{P,Q\}$ with $P\defeq C\ell_{j_f}^{-2-a-\alpha}(|u|^{1+\alpha}+\eta'^{1+\alpha})$ and, by \Cref{rot:h1s:E} for a slice leg, $Q\defeq b_0+b_1|u|+b_2$ with $b_0=C(\eta'/\ell_f)^{1/2^a}\ell_f^{-1-a}$, $b_1=C(\eta'/\ell_f)^{1/2^{a+1}}\ell_f^{-2-a}$, $b_2=C\eta'\ell_f^{-2-a}$ ($\ell_f\defeq\ell_{j_f}$). Since $(1+\alpha)(1-s)=1$, $\min\{P,Q\}\le P^{1-s}Q^s\le C\gamma_f\ell_f^{-2-a}(|u|+\eta')(1+|u|/\ell_f)^s$ with $\gamma_f\defeq(\eta'/\ell_f)^{s/2^{a+1}}$ (the $\ell_f$-exponent is $-(2+a+\alpha)(1-s)-(1+a)s=-2-a$); then $B_\sigma\le C\gamma_f\ell_f^{-2-a}e^{\dots}\sum_u\eta'^3|\kappa_\sigma(u)|(|u|+\eta')$, and the first-order sum of a degree-$(-1)$ kernel is $\le C$ per $\sigma$, hence $\le C(1+r)$ after the $\sigma$-sum: one logarithm. External leg: $b_0=C\eta^{1/2}\|E\|$, $b_1=C\eta^\kappa\|E\|$, $b_2=C\eta'\|E\|$, $\gamma_f=\eta^{\kappa s}\le5\eta'^{\kappa s}$. (iii) Slice $l$: $|T_xT_wl(y,z)|\le\ell_l^{-3-a-2\alpha}|u_1|^{1+\alpha}|u_2|^{1+\alpha}$ (mixed second-order lattice remainder, from \eqref{eq:h1s:B} at order $a+4$) gives, summed, $\ell_l^{-3-a}$ per flow. Difference: $P=\ell_l^{-3-a-2\alpha}|u_1|^{1+\alpha}|u_2|^{1+\alpha}$ and $Q\le C(\eta'/\ell_l)^{1/8}\ell_l^{-1-a}(1+|u_1|/\ell_l)(1+|u_2|/\ell_l)$, since the differenced factors $l,\nabla l,\nabla\nabla l$ carry $n\le a+2\le3$ derivatives (\Cref{rot:h1s:P}, exponent $\frac18$) plus the displacements; $P^{1-s}Q^s\le C(\eta'/\ell_l)^{s/8}\ell_l^{-3-a}|u_1||u_2|(\dots)^s$ as $(3+a+2\alpha)(1-s)+(1+a)s=3+a$; two first-order sums give $(1+r)^2$. Coarse $l=\Gamma^F=H^FB^FH^{F*}$ (\Cref{rot:lem:Gam}): $T_xT_w$ acts on the two $H$-factors separately, each $C^{1,\alpha}$ in its first variable with $|T_xH^X(\cdot,p)(y)|\le C|y-x|^{1+\alpha}e^{-\delta|x-p|}$, and $|B^F(p,q)|\le Ce^{-\delta|p-q|}$, so $|T_xT_w\Gamma^F(y,z)|\le C|u_1|^{1+\alpha}|u_2|^{1+\alpha}e^{-\delta|x-w|/2}$; the difference uses $|\Gamma^X-\Gamma^Y\circ\tau|\le C\eta^{1/2}$ and \Cref{rot:h1s:E} for each $H$-factor, $Q\le C\eta^\kappa(1+|u_1|)(1+|u_2|)$, whence $P^{1-s}Q^s\le C\eta^{\kappa s}|u_1||u_2|$. (iv) First-order forms: $P=\ell_f^{-2-a}|u|$ (slice; $C\|E\||u|$ external), $Q$ as above with $n\le a+1$ in the gradient; $P^{1/2}Q^{1/2}\le C(\eta'/\ell_f)^{1/2^{a+2}}\ell_f^{-1-a}|u|^{1/2}(\dots)$ resp. $C\eta^{1/4}\|E\||u|^{1/2}(\dots)$, and the first-order sum converges per $\sigma$ with the $(1+r)$ after the $\sigma$-sum. For the brackets of (R2b) the switched factor $\Delta[g(t)-g(x)]$ is $\le|t-x|$ times the slice law for $\nabla g$ plus $2c_\tau\eta'\sup|\nabla g^Y|$: the degree-one term with its gain and the explicit-$\eta'$ term (degree zero, one logarithm). When the differenced factor is a compound $(U,f)$ at an outer vertex $t\in\{x,w\}$, its lattice difference is the compound with the differenced leg: at $t=w$ directly, and at the unit vertex $t=x$ by the translation invariance of $\kappa_\sigma$, since shifting the summation variable gives $K_\sigma(x+\eta'e_\mu,w)-K_\sigma(x,w)=\eta'\sum_y\kappa_\sigma(x-y)\,T_x(\partial_\mu f)(y)$ exactly, the lattice differences commuting (and likewise on $\Lambda_Y$ with the step $\eta'Re_\mu$). So the difference of the bracket is a sum of $|t-x|/\eta'$ terms $\eta'\,[K^X_\sigma-K^Y_\sigma]$ with the leg $\partial_\mu f$, $a+1\le2$, each bounded by (ii) (\Cref{rot:h1s:E} for the slice leg with $a+1$, \Cref{rot:h1s:P} at $n\le3$; the unit lines keep $n\le2$): per flow $|t-x|\,\ell_{j_f}^{-3-a}$, gain $\eta'^{\theta_U}$ with $\theta_U=\min\{\theta,s/2^{a+2}\}\ge\min\{\theta,s/8\}$, one logarithm; the $Y$-displacement term is explicit-$\eta'$ as before.
\end{proof}
